\chapter{Mathematical Formulas \& Vector Identities}

\section{Differential Operators in Curvilinear Coordinates}
Let $\psi$ denote a smooth scalar field and $\vec{A} = A_1 \hat{e}_1 + A_2 \hat{e}_2 + A_3 \hat{e}_3$ denote a continuously differentiable vector field.

\subsection{Cartesian Coordinates $(x, y, z)$}
Metric scale factors: $h_x = 1, \, h_y = 1, \, h_z = 1$.
\begin{align}
\nabla \psi &= \frac{\partial \psi}{\partial x}\hat{i} + \frac{\partial \psi}{\partial y}\hat{j} + \frac{\partial \psi}{\partial z}\hat{k} \\
\nabla \cdot \vec{A} &= \frac{\partial A_x}{\partial x} + \frac{\partial A_y}{\partial y} + \frac{\partial A_z}{\partial z} \\
\nabla \times \vec{A} &= \left(\frac{\partial A_z}{\partial y} - \frac{\partial A_y}{\partial z}\right)\hat{i} + \left(\frac{\partial A_x}{\partial z} - \frac{\partial A_z}{\partial x}\right)\hat{j} + \left(\frac{\partial A_y}{\partial x} - \frac{\partial A_x}{\partial y}\right)\hat{k} \\
\nabla^2 \psi &= \frac{\partial^2 \psi}{\partial x^2} + \frac{\partial^2 \psi}{\partial y^2} + \frac{\partial^2 \psi}{\partial z^2}
\end{align}

\subsection{Cylindrical Coordinates $(\rho, \phi, z)$}
Metric scale factors: $h_\rho = 1, \, h_\phi = \rho, \, h_z = 1$.
\begin{align}
\nabla \psi &= \frac{\partial \psi}{\partial \rho}\hat{\rho} + \frac{1}{\rho}\frac{\partial \psi}{\partial \phi}\hat{\phi} + \frac{\partial \psi}{\partial z}\hat{k} \\
\nabla \cdot \vec{A} &= \frac{1}{\rho}\frac{\partial}{\partial \rho}(\rho A_\rho) + \frac{1}{\rho}\frac{\partial A_\phi}{\partial \phi} + \frac{\partial A_z}{\partial z} \\
\nabla \times \vec{A} &= \left(\frac{1}{\rho}\frac{\partial A_z}{\partial \phi} - \frac{\partial A_\phi}{\partial z}\right)\hat{\rho} + \left(\frac{\partial A_\rho}{\partial z} - \frac{\partial A_z}{\partial \rho}\right)\hat{\phi} + \frac{1}{\rho}\left(\frac{\partial (\rho A_\phi)}{\partial \rho} - \frac{\partial A_\rho}{\partial \phi}\right)\hat{k} \\
\nabla^2 \psi &= \frac{1}{\rho}\frac{\partial}{\partial \rho}\left(\rho\frac{\partial \psi}{\partial \rho}\right) + \frac{1}{\rho^2}\frac{\partial^2 \psi}{\partial \phi^2} + \frac{\partial^2 \psi}{\partial z^2}
\end{align}

\subsection{Spherical Coordinates $(r, \theta, \phi)$}
Metric scale factors: $h_r = 1, \, h_\theta = r, \, h_\phi = r\sin\theta$.
\begin{align}
\nabla \psi &= \frac{\partial \psi}{\partial r}\hat{r} + \frac{1}{r}\frac{\partial \psi}{\partial \theta}\hat{\theta} + \frac{1}{r\sin\theta}\frac{\partial \psi}{\partial \phi}\hat{\phi} \\
\nabla \cdot \vec{A} &= \frac{1}{r^2}\frac{\partial}{\partial r}(r^2 A_r) + \frac{1}{r\sin\theta}\frac{\partial}{\partial \theta}(\sin\theta A_\theta) + \frac{1}{r\sin\theta}\frac{\partial A_\phi}{\partial \phi} \\
\nabla \times \vec{A} &= \frac{1}{r\sin\theta}\left(\frac{\partial (\sin\theta A_\phi)}{\partial \theta} - \frac{\partial A_\theta}{\partial \phi}\right)\hat{r} \nonumber\\
&\quad + \frac{1}{r}\left(\frac{1}{\sin\theta}\frac{\partial A_r}{\partial \phi} - \frac{\partial (r A_\phi)}{\partial r}\right)\hat{\theta} + \frac{1}{r}\left(\frac{\partial (r A_\theta)}{\partial r} - \frac{\partial A_r}{\partial \theta}\right)\hat{\phi} \\
\nabla^2 \psi &= \frac{1}{r^2}\frac{\partial}{\partial r}\left(r^2\frac{\partial \psi}{\partial r}\right) + \frac{1}{r^2\sin\theta}\frac{\partial}{\partial \theta}\left(\sin\theta\frac{\partial \psi}{\partial \theta}\right) + \frac{1}{r^2\sin^2\theta}\frac{\partial^2 \psi}{\partial \phi^2}
\end{align}

\section{Fundamental Vector and Differential Identities}
\begin{align}
\vec{A} \cdot (\vec{B} \times \vec{C}) &= \vec{B} \cdot (\vec{C} \times \vec{A}) = \vec{C} \cdot (\vec{A} \times \vec{B}) \\
\vec{A} \times (\vec{B} \times \vec{C}) &= (\vec{A} \cdot \vec{C})\vec{B} - (\vec{A} \cdot \vec{B})\vec{C} \\
\nabla \cdot (\nabla \times \vec{A}) &= 0 \quad (\text{Divergence of a curl is zero}) \\
\nabla \times (\nabla \psi) &= \vec{0} \quad (\text{Curl of a gradient is zero}) \\
\nabla \times (\nabla \times \vec{A}) &= \nabla(\nabla \cdot \vec{A}) - \nabla^2 \vec{A} \\
\nabla (\vec{A} \cdot \vec{B}) &= \vec{A} \times (\nabla \times \vec{B}) + \vec{B} \times (\nabla \times \vec{A}) + (\vec{A} \cdot \nabla)\vec{B} + (\vec{B} \cdot \nabla)\vec{A}
\end{align}
